\documentclass[11pt]{article}
\usepackage[a4paper,margin=26mm]{geometry}
\usepackage{amsmath,amssymb,amsthm,mathrsfs,hyperref}
\setlength{\emergencystretch}{3em}
\newtheorem{theorem}{Theorem}
\newtheorem{lemma}[theorem]{Lemma}
\newtheorem{proposition}[theorem]{Proposition}
\newcommand{\PSH}{\operatorname{PSH}}
\newcommand{\vol}{\operatorname{vol}}
\newcommand{\MA}{\operatorname{MA}}
\newcommand{\EE}{\mathcal E^1_\theta}
\newcommand{\OO}{\mathcal O}
\title{Complete moving systems and actual relative contraction}
\author{Unreviewed research, continuation 208}
\date{6 October 2026}
\begin{document}
\maketitle
\begin{abstract}
For an actual big rational bundle on a smooth projective variety
over an algebraically closed trivially valued characteristic-zero
field, we reconstruct all lower moving products of complete systems.
Full time lattices, with their entire pulled horizontal fixed factor
removed before principalization, give a common generated generic
bundle for every fixed finite list of obstacles. Their auxiliary
energies and original pushed measures have the claimed limits.
A retained missing-rank section trace and a finite nef energy
polynomial prove quadratic contraction relative to the intrinsic
minimal reference. This gives an actual finite-energy pole on a
Kodaira divisor and the equivalence of all-test integrability with
finite intrinsic dual energy. The companion outer-energy construction
is fully specified proof data requiring joint independent review.
No unbounded physical weighted reference-defect limit or complete
original equation is asserted here.
\end{abstract}

\section{Actual domain and proof dependencies}
Retain every smooth integral projective $n$-fold $X$, $n\geq1$,
over every algebraically closed trivially valued characteristic-zero
field, and every actual big rational bundle $L$.
Set $K=X^{\mathrm{an}}$, $\theta=c_1(L)$, $V=\vol(L)>0$,
and $o$ equal to the generic trivial valuation.
All supported semivaluations and arbitrary directed nets are retained.
The genuine cone is the literal all-point single decreasing FS-net
cone. Its actual scalar closure and envelope theorem is the checked
scalar result of round 160. The checked unconditional finite result
of round 159 supplies bounded nef intersection calculus, positive
raw measures and the full section-lattice intercept.
Neither paper's conditional whole S--M clauses are imported.

The companion full proof 207 constructs on this actual cone
\[
 R=P_\theta(0),\quad E(u)=\inf_{f\in C(K),\ f\geq u}\mathscr E(f),
 \quad \EE=\{E> -\infty\},\quad B(u)=-E(u)\ (u(o)=0),
                                                        \tag{1}
\]
including decreasing energy continuity, nonpositive canonical model
majorants, and a positive target-independent raw-$V$ derivative $M$.
It proves $M(P_\theta(f_D))=\mu_D$, where $\mu_D$ is the finite
positive-product derivative. We use that written companion proof,
not an unexplained whole-energy or whole-operator axiom.
The new whole proofs 207 and 208 require joint independent review.

Ordinary algebraic positive products are used only on smooth
projective characteristic-zero determinations. We use their
directed-supremum definition, nef monotonicity, compactness and
Fujita equality $\vol(N)=\langle N^d\rangle$ from BFJ
\cite[Sections 2--3]{BFJ}. Those complete pertinent units were
read, and their precise algebraic domain was independently checked
in round 161. Actual ample generation and normal birational section
descent retain their verified algebraic scopes. No section-ring
finite generation or field-countability premise is added.

\section{Cofinality of all lower complete moving products}\label{cofinal}
Let $Z$ be smooth integral projective of dimension $d$ and $N$ an
actual big rational bundle. In each cleared nonempty degree resolve
the base ideal of the \emph{full} system:
\[
 p_m^*(mN)=M_m+G_m,\qquad A_m=M_m/m,
 \quad M_m\text{ generated},\quad G_m\geq0.             \tag{2}
\]
All comparisons use common smooth dominations.
The products of sections give $A_{m\ell}\geq A_m$ in effective
divisor order. Degrees are directed by divisibility; a factorial
subnet eventually clears every fixed denominator and degree.
Every sufficiently large cleared degree is nonempty by an actual
Kodaira decomposition and eventual generation of its ample factor.

\begin{theorem}\label{products}
For every $1\leq p\leq d$,
\[
                       A_m^p\longrightarrow\langle N^p\rangle
                                                        \tag{3}
\]
in the weak numerical cycle topology. In particular
$A_m^d\to\vol(N)$. For every fixed signed Cartier test $H$ on
a fixed determination,
\[
 A_m^{d-1}H\longrightarrow\langle N^{d-1}\rangle H.     \tag{4}
\]
\end{theorem}
\begin{proof}
On one fixed determination take rational nef competitors
$\beta_1,\ldots,\beta_p$ with $N-\beta_i$ pseudoeffective.
Choose an actual ample rational $Q_0$ with $N-Q_0$ big there.
For rational $\epsilon>0$ put
$Q_i=(1-\epsilon)\beta_i+\epsilon Q_0$.
These classes are ample and $N-Q_i$ is big, because it is the sum
of $(1-\epsilon)(N-\beta_i)$ and the big class
$\epsilon(N-Q_0)$. Choose actual rational representatives.
In one sufficiently large cleared degree $q$, all $qQ_i$ are
generated and all $q(N-Q_i)$ have nonzero sections with divisors
$R_i$. Multiplication by these sections puts their generated
systems inside the \emph{full} $|qN|$.
Normal birational descent identifies this latter section space
with the original one. The pulled base ideal contains
$\OO(-R_i)$, so $G_q\leq R_i$ and $A_q\geq Q_i$.

Nef monotonicity now gives $A_q^p\geq Q_1\cdots Q_p$.
Let $\epsilon\downarrow0$ on the fixed determination.
Thus any upper bound for all $A_m^p$ is an upper bound for every
defining competitor product. Conversely each $A_m$ is itself an
allowed nef competitor under $N$. Their least upper bound is
therefore $\langle N^p\rangle$.
Real nef competitors follow by approximation from the ample
$Q_i$ before this limit. BFJ's directed compactness argument
\cite[Lemmas 2.6--2.7]{BFJ} makes the directed net converge to
this least upper bound: every cluster is an upper bound for every
fixed stage and lies below every upper bound; compactness supplies
clusters and uniqueness identifies the whole net.

For the signed $H$, choose two fixed ample rational numerical
classes $H_1,H_2$ on its determination with $H=H_1-H_2$.
Pair (3) with them separately, then subtract.
The test and its decomposition are fixed before the degree limit.
There is no effective-order comparison against a signed test.
The top assertion follows from the algebraic Fujita equality.
\end{proof}

\section{Full time lattices and one common generic moving bundle}
Fix a finite list of rational model obstacles $f_D\leq0$,
$f_D(o)=0$. On a smooth equivariant determination write
$F=\sum_E b_EE$ and $\alpha_D(c)=\rho_X^*L+D+cF$.
The full relative lattice has an adapted basis with weights
$-mC\leq w_{mj}\leq0$ for fixed $C$.
For one integer $c>C+1$, the \emph{full absolute} section space
of $m\alpha_D(c)$ has the basis
\[
 t^{-(w_{mj}+mc)+r}s_{mj},\qquad 0\leq r\leq w_{mj}+mc.
                                                        \tag{5}
\]
This is the lattice dimension formula, including every time weight.
At nonzero time it recovers every $s_{mj}$, and the endpoint at
infinity also recovers the entire generic system.

Resolve the full $|mL|$ once and write the actual generic identity
\[
 p_m^*L=A_m+T_m,\qquad mA_m\text{ generated},\quad mT_m\geq0.
                                                        \tag{6}
\]
The weights in the finite list may use different adapted bases of
the \emph{same} whole space $H^0(X,mL)$.
On a common equivariant model over $X_m\times\mathbb P^1$,
every section in (5) has the full actual factor $s_{mT_m}$.
Remove its entire pullback, including any vertical components,
before principalizing each residual ideal.
The residual ideals are unit ideals away from zero because the
remaining generic full systems generate $mA_m$.
Their principalized divisors $G^{\rm vert}_{m,D}$ are consequently
effective and supported at zero.
The resulting full moving bundle is
\[
 C_{m,D}=\rho_X^*A_m+D+cF-G^{\rm vert}_{m,D}/m.          \tag{7}
\]
It is globally generated, hence absolutely nef.
Taking a common domination for the finite list preserves nefness
and leaves the same actual generic $A_m$ in every slot.

Let $\psi_{m,D}$ be its bounded relatively generated model scalar
on $A_m$. The weighted full generic basis gives $-C\leq\psi_{m,D}\leq0$.
The actual pulled section identities are
\[
 \phi_{m,D}=q_{T_m}+\psi_{m,D},\qquad \phi_{m,0}=q_{T_m}.
                                                        \tag{8}
\]
Here $q_{T_m}$ is the canonical logarithmic norm of its section.
These physical functions can have poles; no original finite-energy
membership is assigned to them.
The full filtration gives the all-point minorant
$\phi_{m,D}\leq P_\theta(f_D)$ on every lift.

\begin{theorem}[Whole finite-list moving signature]\label{moving}
For the fixed finite list in the actual big $L$,
\[
\begin{split}
 V_m=A_m^n&\longrightarrow V,\\
 e_{A_m}(\psi_{m,D})
   =\frac{C_{m,D}^{n+1}}{n+1}-cV_m
      &\longrightarrow E(P_\theta(f_D)),\\
 p_{m*}\MA_{A_m,\mathrm{raw}}(\psi_{m,D})
      &\longrightarrow\mu_D=M(P_\theta(f_D))
\end{split}                                             \tag{9}
\]
weakly in the last line. Every coefficient retains its central
multiplicity $b_E$. Generic values $\psi_{m,D}(o)$ tend to zero.
\end{theorem}
\begin{proof}
The construction (7) is precisely the moving part of the entire
absolute system (5), not of a selected subseries.
Theorem~\ref{products}, applied to $L$ and to each $\alpha_D(c)$,
gives $V_m\to V$, $C_{m,D}^{n+1}\to\vol(\alpha_D(c))$,
and $C_{m,D}^nH\to\langle\alpha_D(c)^n\rangle H$
for each fixed signed original vertical Cartier test $H$.
The full section intercept and companion 207's envelope identity
give the energy limit.
The finite nef raw measure has coefficient
$b_E(C_{m,D}|_E)^n$ at the associated valuation; its integral
against the pulled model test represented by $H$ is $C_{m,D}^nH$.
The signed-test limit identifies it with $\mu_D(H)$.
Uniform original model density and raw masses $V_m\to V$
give weak convergence against every $h\in C(K)$.
This uses lower moving products, not only top-volume convergence.
Companion 207's touching argument identifies the final measure with
$M(P_\theta(f_D))$.

For the generic limit, actual bounded-obstacle section density
is enough; we give its proof next. The construction is simultaneous
for any fixed finite list by denominator clearing and common
principalization. No moving unbounded amplitude is involved.
\end{proof}

\section{Actual bounded section density off a Kodaira divisor}
Choose an actual Kodaira identity $L=aA+E_0$, $A$ ample,
$a>0$ rational, $E_0\geq0$. Put $W=q_{E_0}$, $Z=E_0^{\mathrm{an}}$,
and $U=K\setminus Z$. Tensoring its section with a generated
power of $aA$ makes $W$ a genuine FS scalar on $L$.
For bounded continuous $f$, use its full weights
$W_m(f;s)=\inf_K(mf-\log|s|)$ and an adapted full basis to define
\[
 \phi_{m,f}=\frac1m\max_j
                 (\log|s_{mj}|+\lfloor W_m(f;s_{mj})\rfloor).
                                                        \tag{10}
\]
The definition is basis independent by the ultrametric inequality
on each filtration space. It increases on the divisibility-directed
degrees and lies below the genuine envelope $u_f=P_\theta(f)$.

\begin{lemma}\label{density}
On $U$, $\sup_m\phi_{m,f}=u_f$, the latter is finite continuous,
and the cofinal stages converge uniformly on every compact subset.
In particular $q_{T_m}\to R$ locally uniformly on $U$,
$\psi_{m,D}\to P_\theta(f_D)-R$ locally uniformly there,
and $\psi_{m,D}(o)\to0$.
\end{lemma}
\begin{proof}
Choose an integer $C\geq\|f\|_\infty$.
The genuine minorant $W-C$ makes $u_f$ finite on $U$.
Take a defining FS net $h_i\downarrow u_f$ on actual $B_i$.
For $\epsilon>0$ compact strict margins give eventually
$h_i\leq f+\epsilon$ on all of $K$.
Fix rational $t\in(0,1)$ before this tail. The actual bundle
\[
                   L-(1-t)B_i-t(L-aA)
\]
is eventually ample because its class tends to $ta[A]$.
Its weight-zero generated system is identically zero at every
original semivaluation. Tensor it with the FS stages to produce
the actual $L$-FS minorant
$(1-t)(h_i-\epsilon)+t(W-C)$ of $f$, with exponents cleared.

Every actual FS minorant $g$ of $f$ is below $\sup_m\phi_{m,f}$:
raise its finitely many sections to powers; their full feasible
weights scale and the integer floor loses less than $1/m$.
Apply this to the preceding tensor minorant, take the defining
tail at fixed $x\in U$, then $t\downarrow0$ and
$\epsilon\downarrow0$. This proves density at $x$.
All sufficiently divisible full systems have no base point on $U$
by the Kodaira identity. Their stages are finite continuous there.
The increasing supremum is lsc and the genuine envelope is usc,
so it is continuous. On a fixed compact subset the continuous
gaps decrease to zero; compact strict margins give directed Dini
convergence. Apply this to $f=0$ and to each $f_D$, then subtract
the finite values in (8). Finally $o\in U$, and the normalized
envelope and reference values there are zero.
\end{proof}

\section{The exact missing-rank trace inequality}\label{trace}
Fix $m$ and a bounded nonpositive relatively nef model scalar
$\psi$ on the actual generated $A_m$. Its full section trace
is its finite intersection energy $e_{A_m}(\psi)$.
Indeed choose a sufficiently large time twist: vertical curves
are controlled by relative nefness, and horizontal curves by
the effective inequality $D_\psi+cF\geq0$ and the nef generic
pullback. The compactified bundle is nef. Its ordinary volume
is its top intersection; the full lattice intercept identifies
the trace. Finite raw measures and mixed energy formulas follow
from G0 by adding an actual common ample perturbation on each
fixed model and taking the finite polynomial limit. This works
also when $V_m=0$, without dividing by its mass.

In every cleared degree $\ell$, multiplication by the actual
fixed section gives an injection
\[
 H^0(X_m,\ell A_m)\hookrightarrow H^0(X,\ell L).       \tag{11}
\]
Normal proper birational descent is used after multiplication.
Suppose $f\in C(K)$ majorizes $q_{T_m}+\psi$ on every lift.
Every auxiliary feasible section weight remains feasible for
the original $f$ by the literal product section identity; at
a zero of the fixed factor the original section inequality
still holds. No common-pole subtraction is needed.
All remaining original weights are at least $\ell\min_K f$
because original canonical section norms are at most one.
Extend the injected filtration basis to a full adapted basis.
The quotient weights retain this lower bound. With
$r_\ell=h^0(X,\ell L)$ and $d_\ell=h^0(X_m,\ell A_m)$,
the integer traces satisfy, up to an error bounded by the ranks,
\[
 \operatorname{Tr}_{\ell,f}\geq\operatorname{Tr}_{\ell,\psi}
                   +(r_\ell-d_\ell)\ell\min_K f.
\]
First send $\ell\to\infty$ at fixed $m,f,\psi$.
Full section growth and the identified nef trace give
\[
 \mathscr E(f)\geq e_{A_m}(\psi)+(\min_K f)(V-V_m).     \tag{12}
\]
The last term may be negative and cannot be dropped at finite $m$.
Only a subsequent moving limit at this fixed continuous $f$
removes it. Taking an infimum over $f$ earlier is not justified.

\section{Quadratic relative contraction on the whole actual space}
For a bounded nonpositive nef scalar $\psi$ on $A_m$, set
\[
 a_k=\int(-\psi)\,\MA_{A_m}(\psi^k,0^{n-k})\geq0,
 \quad B_m(\psi)=\frac1{n+1}\sum_{k=0}^n a_k.
\]
All mixed raw measures have mass $V_m$; $a_0=-\psi(o)V_m$.
Finite intersection multilinearity gives for $0\leq r\leq1$
\[
 B_m(r\psi)=\frac r{n+1}\sum_{j=0}^n\sum_{k=0}^j
       \binom{j}{k}r^k(1-r)^{j-k}a_k
 \leq r a_0+C_n r^2 B_m(\psi),\quad C_n=2^{n+1}.       \tag{13}
\]
The $k=0$ part is at most $ra_0$; every other summand has a
second factor $r$. Bounding the finite binomial sums proves
the displayed dimension constant. The generic linear term is
retained until the limit.

\begin{theorem}\label{contraction}
Every normalized actual $u\in\EE$ satisfies
\[
           B((1-r)R+ru)\leq C_n r^2 B(u)
                         \quad(0\leq r\leq1).          \tag{14}
\]
\end{theorem}
\begin{proof}
First $u=P_\theta(f_D)$ with $f_D\leq0$, $f_D(o)=0$.
The physical stages give
$q_{T_m}+r\psi_{m,D}\leq(1-r)R+ru$ on every lift.
For each fixed continuous majorant $f$ of the right side apply
(12) and (13):
\[
 \mathscr E(f)\geq
       r\psi_{m,D}(o)V_m-C_n r^2 B_m(\psi_{m,D})
                         +(\min_K f)(V-V_m).
\]
Take the moving limit (9) and Lemma~\ref{density} first.
The retained generic and rank errors vanish, leaving
$\mathscr E(f)\geq-C_n r^2 B(u)$.
Only now take the infimum over these majorants.
For general $u$, companion 207's nonpositive canonical model
envelopes decrease to $u$ with decreasing energies; their
contracted genuine functions also decrease to the desired one.
Decreasing energy continuity passes the inequality.
Rational $r$ suffices for section tensor identities.
Finite concavity gives continuity along the interior of the
genuine segment, hence all real $r$ there. At $r=0$ the reference
has zero cost; at $r=1$ the inequality is immediate.
\end{proof}

\section{A genuine finite-energy pole and intrinsic target conditions}
\begin{proposition}\label{pole}
There is a normalized actual $w_Z\in\EE$ with $w_Z=-\infty$ on
the full actual Kodaira divisor $Z$.
\end{proposition}
\begin{proof}
Let $N_j=16^j$, $r_j=8^{-j}$, and
$a_j=\max(W,R-N_j)$, a normalized genuine scalar of cost at most
$N_jV$. Put $b_j=(1-r_j)R+r_j a_j$.
Theorem~\ref{contraction} gives $B(b_j)\leq C_nV4^{-j}$.
The genuine partial sums
\[
               w_N=\sum_{j=1}^N2^{-j}b_j+2^{-N}R
\]
decrease, are normalized, and have cost at most
$C_nV\sum_{j\geq1}8^{-j}=C_nV/7$ by concavity.
Actual decreasing closure and energy continuity give an actual
finite-energy endpoint $w_Z$. Where $R$ is finite on $Z$,
$R-b_j=r_jN_j=2^j$, and each weighted term contributes one
to $R-w_N$; hence the endpoint is $-\infty$ there.
Where $R=-\infty$, every normalized genuine function is already
$-\infty$. Generic value zero and the cost bound are retained.
\end{proof}

Let $\nu$ be a positive original Radon measure of raw mass $V$
with $\nu\{R=-\infty\}=0$. Define, only on the complement of
that reference pole set, for normalized actual $u\in\EE$,
\[
 I_\nu(u)=\int(R-u)\,d\nu\in[0,\infty],\qquad
 D_\nu=\sup_{u\in\EE,\,u(o)=0}\{I_\nu(u)-B(u)\}.       \tag{15}
\]
No absolute reference integrability is imposed.

\begin{theorem}\label{targets}
$D_\nu<\infty$ if and only if every actual finite-energy relative
test is $\nu$-integrable. For every normalized $u\in\EE$ and
$0<\epsilon\leq C_n$,
\[
 I_\nu(u)\leq\epsilon B(u)+C_nD_\nu/\epsilon.           \tag{16}
\]
Such a target misses $Z$, and any maximizing sequence for
$I_\nu-B$ has uniformly bounded normalized actual energy cost.
\end{theorem}
\begin{proof}
Finite $D_\nu$ directly bounds every $I_\nu$ by $B+D_\nu$.
Conversely suppose every test is integrable.
Their integrals are uniformly bounded on the unit cost sublevel:
otherwise choose $u_j$ there with $I_\nu(u_j)\geq4^j$.
The genuine decreasing partial sums
$\sum_{j=1}^N2^{-j}u_j+2^{-N}R$ have cost at most one.
Their actual endpoint is finite energy; its relative gap is the
countable increasing sum of the gaps off the reference poles.
Ordinary countable monotone convergence makes its integral infinite,
a contradiction. Denote the unit bound by $C_\nu$.

If $B(u)\geq1/C_n$, contract with
$r=(C_nB(u))^{-1/2}$. Its cost is at most one by (14), so
$I_\nu(u)\leq C_\nu\sqrt{C_nB(u)}$.
For smaller costs use the unit bound directly.
The supremum of this growth bound minus $B$ is finite, proving
the converse. The reference gives $D_\nu\geq0$.

For any $0<r\leq1$, (14)--(15) give
$rI_\nu(u)\leq C_n r^2B(u)+D_\nu$.
Set $r=\epsilon/C_n$ to obtain (16).
Taking $\epsilon<1$ gives coercivity of $I_\nu-B$ and bounds
the cost of any near-maximizers; no attainment is inferred from
this bound alone. Finally integrability of the pole in
Proposition~\ref{pole}, together with reference-pole nullness,
forces $\nu(Z)=0$.
\end{proof}

\section{Scope and outstanding joint review}
The actual scalar theorem 160, unconditional finite theorem 159,
and full companion 207 proof are the exact dependencies.
The preceding reconstruction supplies all lower complete moving
products, the fixed finite-list energy and weak-output signature,
the exact missing-rank trace, whole quadratic relative contraction,
an actual finite-energy Kodaira pole, and intrinsic coercivity.
It does not assume the finite-list moving or contraction axioms
that appeared as conditional signatures in earlier papers.

The joint new construction is unreviewed. Neither weak moving
convergence nor the pole proposition alone identifies a fixed
unbounded original test at every physical auxiliary output.
The discrepancy $R-q_{T_m}$ remains in that stronger assertion.
Actual inverse attainment, all genuine mixed and weighted limits,
and the full original equation still require their complete proof
chains and independent review. Original P16.5 stays open;
no original or supporting-paper acceptance count changes.
No new primary retrieval or recursive proof of all background is
claimed. The existing user editor remains unchanged.
\begin{thebibliography}{9}
\small
\bibitem{BFJ} S. Boucksom, C. Favre and M. Jonsson,
\emph{Differentiability of volumes of divisors and a problem of Teissier},
\href{https://arxiv.org/html/math/0608260v2}{arXiv:math/0608260v2}.
\bibitem{LM} R. Lazarsfeld and M. Musta\c t\u a,
\emph{Convex Bodies Associated to Linear Series}, Ann. Sci. ENS,
42 (2009), 783--835.
\bibitem{Companion} Research continuation 207,
\emph{Actual outer section energy and its full-volume derivative},
\href{https://github.com/MingchenXia/nonarchimedean-pluripotential-theory/blob/a61bb373885ed4a270f528a5eae0d51799931df8/research-checkpoints/continuation172/root-checks/P16.5-actual-outer-energy-and-full-volume-derivative-candidate-20261006-continuation207.tex}{Frozen unreviewed source}, 6 October 2026.
\end{thebibliography}
\end{document}
